\chapter{Strings}


\section{Permutaciones de string}
\begin{lstlisting}[language=C++, caption=Next\_permutation]

int main() {
    string s = "ABA";
    
    IMPORTANTE: Ordenar la cadena primero
    sort(s.begin(), s.end()); 
    
    El bucle do-while asegura que se procese la 1 permutacion
    do {
        cout << s << endl;
    } while (next_permutation(s.begin(), s.end()));
    
    return 0;
}
\end{lstlisting}
\textbf{Complejidad:} Para una cadena de longitud $n$ es $n!$.

\section{Palindromic Tree (Eertree)}

\subsection{¿Qué es?}

El \textbf{Palindromic Tree} (también llamado \textit{Eertree}) es una estructura de datos que representa \textbf{todos los substrings palíndromos distintos} de un string $s$ de longitud $n$. Tiene a lo sumo $n+2$ nodos: uno por cada palíndromo distinto, más dos raíces imaginarias.

\bigskip

\textbf{¿Cuándo usarlo?}
\begin{itemize}
    \item Contar la cantidad de substrings palíndromos distintos.
    \item Contar cuántas veces aparece cada palíndromo como substring.
    \item Problemas que requieren explotar la estructura de palíndromos (DP sobre palíndromos, etc.).
\end{itemize}

\subsection{Estructura interna}

Cada nodo del árbol representa un palíndromo distinto que aparece en $s$.

\begin{itemize}
    \item \texttt{len}: longitud del palíndromo que representa este nodo.
    \item \texttt{link}: \textit{suffix link} — apunta al nodo del palíndromo sufijo propio más largo de este nodo.
    \item \texttt{to[c]}: transición: si al nodo $v$ le agregamos el carácter $c$ a ambos lados, llegamos al nodo \texttt{to[c]}.
    \item \texttt{cnt}: cantidad de veces que este palíndromo aparece como sufijo al procesar el string (se propaga al final para obtener ocurrencias totales).
\end{itemize}

\textbf{Raíces especiales:}
\begin{itemize}
    \item Nodo 0: raíz imaginaria con \texttt{len = -1} (palíndromo de longitud $-1$, sirve para el caso base).
    \item Nodo 1: raíz imaginaria con \texttt{len = 0} (palíndromo vacío).
\end{itemize}

\subsection{Implementación}

\begin{lstlisting}[style=cppstyle, caption={Palindromic Tree (Eertree)}]
struct palindromic_tree {
    static const int SIGMA = 26;
    struct Node {
        int len, link, to[SIGMA];
        ll cnt;
        Node(int len, int link = 0, ll cnt = 1)
            : len(len), link(link), cnt(cnt) {
            memset(to, 0, sizeof(to));
        }
    };
    vector<Node> ns;
    int last;

    palindromic_tree() : last(0) {
        ns.pb(Node(-1)); // nodo 0: raiz imaginaria len=-1
        ns.pb(Node(0));  // nodo 1: raiz vacia    len=0
    }

    void add(int i, string& s) {
        int p = last, c = s[i] - 'a';
        // Buscar el sufijo palindromo mas largo al que se puede
        // agregar s[i] a la izquierda
        while (s[i - ns[p].len - 1] != s[i]) p = ns[p].link;
        if (ns[p].to[c]) {
            // El palindromo ya existe: solo actualizamos last
            last = ns[p].to[c];
            ns[last].cnt++;
        } else {
            // Crear un nuevo nodo para el nuevo palindromo
            int q = ns[p].link;
            while (s[i - ns[q].len - 1] != s[i]) q = ns[q].link;
            q = max(1, ns[q].to[c]); // suffix link del nuevo nodo
            last = ns[p].to[c] = (int)ns.size();
            ns.pb(Node(ns[p].len + 2, q, 1));
        }
    }
};
\end{lstlisting}

\subsection{¿Qué tengo que pasarle?}

Se llama \texttt{add(i, s)} para cada posición $i$ del string $s$, de izquierda a derecha.

\textbf{Importante:} el string $s$ debe tener un centinela al inicio (posición 0) que no aparezca en el alfabeto, por ejemplo \texttt{'\$'}. Esto evita que el \texttt{while} salga de rango.

\begin{lstlisting}[style=cppstyle, caption={Uso basico del Palindromic Tree}]
string t; cin >> t;
string s = "$" + t; // centinela en posicion 0

palindromic_tree pt;
for (int i = 1; i < (int)s.size(); i++)
    pt.add(i, s);
\end{lstlisting}

\subsection{¿Qué puedo consultar?}

\subsubsection*{Cantidad de palíndromos distintos}

La cantidad de substrings palíndromos distintos es la cantidad de nodos del árbol menos 2 (las dos raíces imaginarias):

\begin{lstlisting}[style=cppstyle, caption={Contar palindromos distintos}]
int distintos = (int)pt.ns.size() - 2;
\end{lstlisting}

\subsubsection*{Ocurrencias de cada palíndromo}

Después de procesar todo el string, \texttt{ns[v].cnt} solo cuenta las veces que el palíndromo aparece como \textit{sufijo} al momento de ser agregado. Para obtener las ocurrencias totales como substring hay que propagar los \texttt{cnt} por los suffix links, de mayor a menor longitud (los nodos están en orden topológico inverso):

\begin{lstlisting}[style=cppstyle, caption={Propagar cnt para obtener ocurrencias totales}]
// Propagar de atras para adelante (orden topologico)
for (int i = (int)pt.ns.size() - 1; i >= 2; i--)
    pt.ns[pt.ns[i].link].cnt += pt.ns[i].cnt;

// Ahora pt.ns[v].cnt = cantidad total de apariciones del palindromo v
\end{lstlisting}

\subsection{Ejemplo de aplicación}

\textbf{Problema:} dado un string, imprimir la longitud y cantidad de ocurrencias de cada substring palíndromo distinto.

\begin{lstlisting}[style=cppstyle, caption={Imprimir todos los palindromos con su frecuencia}]
string t; cin >> t;
string s = "$" + t;

palindromic_tree pt;
for (int i = 1; i < (int)s.size(); i++)
    pt.add(i, s);

// Propagar ocurrencias
for (int i = (int)pt.ns.size() - 1; i >= 2; i--)
    pt.ns[pt.ns[i].link].cnt += pt.ns[i].cnt;

// Imprimir cada palindromo (nodos 2 en adelante)
for (int i = 2; i < (int)pt.ns.size(); i++)
    cout << "len=" << pt.ns[i].len
         << " ocurrencias=" << pt.ns[i].cnt << "\n";
\end{lstlisting}

\textbf{Ejemplo:} para \texttt{s = "abaab"}:
\begin{itemize}
    \item \texttt{"a"} (len=1): 3 ocurrencias
    \item \texttt{"b"} (len=1): 2 ocurrencias
    \item \texttt{"aba"} (len=3): 1 ocurrencia
    \item \texttt{"aa"} (len=2): 1 ocurrencia
    \item \texttt{"baab"} (len=4): 1 ocurrencia
\end{itemize}

\subsection{Complejidad}

\begin{itemize}
    \item \textbf{Tiempo}: $O(n \cdot \Sigma)$ para construcción, donde $\Sigma$ es el tamaño del alfabeto.
    \item \textbf{Espacio}: $O(n \cdot \Sigma)$ para el árbol.
    \item \textbf{Nodos}: a lo sumo $n + 2$ nodos (uno por cada palíndromo distinto).
\end{itemize}